\section*{Chapitre \ref{ch:b1:ecosystem-dynamics} --- La dynamique des écosystèmes}

\begin{solution}{exo:b1:ecosystem-dynamics:1}
Primaire~: sur un substrat neuf sans sol (une moraine, une coulée de
lave)~; secondaire~: là où le sol et la banque de graines subsistent (un
champ abandonné, un bois brûlé). La secondaire est plus rapide, d'un
facteur dix ou davantage, parce que le sol, ses éléments nutritifs et
ses graines n'ont pas à être fabriqués.
\end{solution}

\begin{solution}{exo:b1:ecosystem-dynamics:2}
Facilitation, inhibition, tolérance. Sur la lave, la facilitation~: les
pionnières (\omterm{def:b1:species-interactions:symbiosis}{lichens}, mousses, fixateurs d'azote) doivent bâtir un sol
avant que quoi que ce soit d'autre puisse s'enraciner.
\end{solution}

\begin{solution}{exo:b1:ecosystem-dynamics:3}
\omterm{def:b1:ecosystem-dynamics:disturbance}{Résistance}~: demeurer inchangé au travers d'une \omterm{def:b1:ecosystem-dynamics:disturbance}{perturbation}~;
\omterm{def:b1:ecosystem-dynamics:disturbance}{résilience}~: revenir vite ensuite. Un récif n'est ni résistant au
réchauffement ni résilient~; une prairie a une faible \omterm{def:b1:ecosystem-dynamics:disturbance}{résistance} à
l'incendie et une forte \omterm{def:b1:ecosystem-dynamics:disturbance}{résilience}.
\end{solution}

\begin{solution}{exo:b1:ecosystem-dynamics:4}
$16^{0.25} = 2$~: \num{24} espèces~; $625^{0.25} = 5$~: \num{60} espèces.
\end{solution}

\begin{solution}{exo:b1:ecosystem-dynamics:5}
La moraine ne contient presque pas d'azote~; les épicéas ne savent pas
le fixer et ne poussent que là où il est fourni. Les aulnes, qui fixent
l'azote et portent la teneur du sol à des centaines de kilogrammes par
hectare, sont donc ce qui rend le site propice à l'épicéa~:
facilitation.
\end{solution}

\begin{solution}{exo:b1:ecosystem-dynamics:6}
$S^* = I_0 P/(I_0 + eP) = 1600/(8 + 10) = 89$ espèces~; renouvellement
$eS^* = 4.4$ espèces par an. Avec $I_0 = 4$~: $800/14 = 57$ espèces,
renouvellement \num{2.9}.
\end{solution}

\begin{solution}{exo:b1:ecosystem-dynamics:7}
Les \omterm{def:b1:ecosystem-organization:trophic}{décomposeurs} ont continué de minéraliser la \omterm{def:b1:ecosystem-dynamics:decomposition}{litière} et l'humus en
ammonium, et les nitrifiants l'ont converti en nitrate, un anion
soluble que le sol ne retient pas~; sans \omterm{def:b1:flowering-plant-organization:organs}{racines} pour le prélever, il a
été lessivé. La repousse a rétabli le prélèvement~: les plantes
nouvelles ont pris le nitrate aussi vite qu'il se formait, et le
ruisseau a retrouvé sa limpidité.
\end{solution}

\begin{solution}{exo:b1:ecosystem-dynamics:8}
$k = \ln 2/1 = \qty{0.69}{yr^{-1}}$~; après trois ans, il reste $e^{-2.08} =
\qty{12.5}{\%}$. Conifère~: $k = -\ln 0.85 = \qty{0.16}{yr^{-1}}$,
demi-vie \num{4.3} ans, \qty{61}{\%} restant après trois ans --- la
\omterm{def:b1:ecosystem-dynamics:decomposition}{litière} s'accumule.
\end{solution}

\begin{solution}{exo:b1:ecosystem-dynamics:9}
Sans pâturage~: de hautes graminées compétitrices ombragent le reste,
peu d'espèces. Pâturage modéré~: les dominantes sont broutées, la
lumière atteint le sol, de petites herbacées coexistent avec les
graminées~: diversité maximale. Pâturage intense~: seules les espèces
prostrées, peu appétentes ou à croissance rapide survivent au piétinement
et au broutage~: de nouveau peu d'espèces.
\end{solution}

\begin{solution}{exo:b1:ecosystem-dynamics:10}
Forêt initiale~: $c\,1000^{0.25} = 5.6c$~; un fragment $10^{0.25}c = 1.78c$,
soit \qty{32}{\%} de la forêt initiale. Dix fragments isolés en
contiendraient, si leurs listes étaient entièrement différentes,
$17.8c$ --- plus que la forêt initiale, ce qui est absurde~: les
fragments partagent la plupart de leurs espèces, donc la somme compte
double. Le total réel se situe entre $1.78c$ (listes identiques) et
environ $5.6c$ (la forêt initiale), moins ce que l'isolement et les
lisières retirent~; les espèces perdues les premières sont celles qui
ont besoin de plus de \qty{10}{km^2} chacune.
\end{solution}

\begin{solution}{exo:b1:ecosystem-dynamics:11}
Un rétablissement exact suggérerait une \omterm{def:b1:ecosystem-organization:ecosystem}{biocénose} fixe, déterminée par
les conditions de l'île~; le rétablissement du nombre avec une liste
différente, et une liste qui ne cesse de changer, est ce que prédit un
équilibre entre immigration et extinction aléatoires --- le nombre est
une propriété de l'aire et de la distance, les identités sont affaire
de hasard.
\end{solution}

\begin{solution}{exo:b1:ecosystem-dynamics:12}
Il n'est pas fermé~: il reçoit la pluie, les poussières, l'azote de
l'air (par fixation et dépôt) et des ions issus de l'altération des
roches, et il perd de l'eau, des ions dissous et des gaz par son
ruisseau et par son air~; mais entrées et sorties sont faibles devant le
cycle interne de la \omterm{def:b1:ecosystem-dynamics:decomposition}{litière} au sol, du sol à la \omterm{def:b1:flowering-plant-organization:organs}{racine}, de la \omterm{def:b1:flowering-plant-organization:organs}{racine} à
la \omterm{def:b1:flowering-plant-organization:organs}{feuille}, si bien que la forêt paraît fermée à la résolution d'un
bilan annuel. La coupe a montré d'où venait cette fermeture apparente~:
du prélèvement par la végétation vivante~; celle-ci ôtée, le même sol et
les mêmes \omterm{def:b1:ecosystem-organization:trophic}{décomposeurs} ont fait fuir à la vallée ses éléments nutritifs
en une saison.
\end{solution}

\begin{solution}{pb:b1:ecosystem-dynamics:1}
\textbf{1.} 10--40 ans~: $140/30 = \qty{4.7}{kg/ha}$ par an~; 40--80~:
$150/40 = \num{3.75}$~; 80--200~: $20/120 = \num{0.17}$. Le maximum est
au stade de l'aulne.
\textbf{2.} La facilitation, pour la transition de l'aulne à l'épicéa~:
l'azote qu'ajoutent les aulnes est ce dont les épicéas ont besoin.
\textbf{3.} La forêt a atteint un régime stationnaire~: les entrées
(fixation, dépôt) équilibrent les pertes (lessivage, dénitrification,
bois exporté) et l'azote circule en interne au lieu de s'accumuler.
\textbf{4.} \qty{150}{kg/ha} retenus en 40 ans contre \num{2000} fixés~:
\qty{7.5}{\%}~; le reste est retenu dans la biomasse et la \omterm{def:b1:ecosystem-dynamics:decomposition}{litière} des
aulnes eux-mêmes, lessivé, ou dénitrifié.
\textbf{5.} Planter des semis d'épicéa sur une moraine de 20 ans avec
et sans aulnes, et avec et sans fertilisation azotée~: si l'épicéa
pousse sans aulnes lorsqu'il est fertilisé, et non sans fertilisation,
l'effet des aulnes est l'azote (facilitation)~; si l'épicéa pousse
également dans toutes les parcelles, simplement plus lentement que
l'aulne, c'est la tolérance.
\textbf{6.} Elle suppose que tous les sites sont partis pareils et ont
suivi le même chemin, ce que le climat, les arrivées de hasard et les
moraines variables du glacier ont pu rompre~; seule une parcelle
permanente prouve la séquence en un lieu donné.
\textbf{7.} De \qty{330}{mm}, soit \qty{40}{\%}~: les arbres ne
transpiraient plus leur part de la pluie.
\textbf{8.} Témoin~: \qty{2}{kg/ha} dans \qty{8200}{m^3}~:
\qty{0.24}{mg/L}. Coupée~: \qty{120}{kg} dans \qty{11500}{m^3}~:
\qty{10.4}{mg/L}, quarante fois plus.
\textbf{9.} Témoin~: $6.5 - 2 = +\qty{4.5}{kg/ha}$ par an, en
accumulation. Coupée~: $6.5 - 120 = -\qty{113.5}{kg/ha}$ par an, en
perte.
\textbf{10.} $120/3850 = \qty{3.1}{\%}$ en un an~; le stock durerait
quelque \num{30} ans à ce rythme --- mais le rythme baisse à mesure que
la repousse revient.
\textbf{11.} Le calcium par dix, le potassium par dix-huit. Le
potassium est retenu surtout dans les \omterm{def:b1:body-plans-tissues:tissue}{tissus} vivants et se lessive
facilement depuis la \omterm{def:b1:ecosystem-dynamics:decomposition}{litière}~; le calcium est aussi fixé sur les argiles
du sol et dans la roche, et libéré plus lentement.
\textbf{12.} De la minéralisation de la \omterm{def:b1:ecosystem-dynamics:decomposition}{litière}, de l'humus et des
rémanents de coupe~: l'ammonification (azote organique en ammonium) par
les \omterm{def:b1:ecosystem-organization:trophic}{décomposeurs}, puis la nitrification (ammonium en nitrate) par les
bactéries nitrifiantes, plus rapides dans un sol découvert, plus chaud
et plus humide.
\textbf{13.} En passant de \num{120} à \num{20} puis à \num{5}~: on
passe au-dessous de l'apport de \qty{6.5}{kg/ha} vers la sixième année~; c'est
le prélèvement d'un nouveau couvert de cerisiers de Pennsylvanie, de
framboisiers et de jeunes arbres qui rétablit la rétention.
\textbf{14.} $\log A$~: 0, 1, 2, 3, 4~; $\log S$~: 1,08, 1,32, 1,58,
1,83, 2,08 --- une droite.
\textbf{15.} $z = (2.08 - 1.08)/4 = 0.25$~; $c = 12$. Îles
intermédiaires~: $12\times 10^{0.25} = 21.3$, $12\times 10^{0.5} = 37.9$, $12\times
10^{0.75} = 67.5$~: toutes à une espèce près.
\textbf{16.} $12\times 250^{0.25} = 12\times 3.98 = 48$ espèces.
\textbf{17.} $38 = 6\times 300/(6 + 300e)$~: $6 + 300e = 47.4$, $e =
\qty{0.138}{yr^{-1}}$~; renouvellement $0.138\times 38 = 5.2$ espèces par an.
\textbf{18.} $3\times 300/(3 + 41.4) = 20$ espèces.
\textbf{19.} $30\times 300/(30 + 41.4) = 126$ espèces. L'extinction sur
une île de \qty{100}{km^2} reste rapide~: même avec un accès libre, les
espèces qui ne peuvent maintenir une \omterm{def:b1:populations:population}{population} sur cette aire
continuent de disparaître~; c'est l'aire, non l'accès, qui fixe le
plafond.
\textbf{20.} Un $z$ plus élevé signifie que le nombre d'espèces chute
plus vite avec l'aire~: les petites îles sont relativement plus pauvres.
Biologiquement, les \omterm{def:b1:populations:population}{populations} des petites îles s'éteignent plus
facilement, ou bien les îles sont si isolées que peu d'espèces les
atteignent (l'isolement élève $z$).
\textbf{21.} $(1000/100\,000)^{0.25} = 0.01^{0.25} = 0.32$~: environ un
tiers.
\textbf{22.} Isolée, chacune en contient $(100/100\,000)^{0.25} = 0.18$~;
si leurs listes étaient indépendantes, elles en contiendraient ensemble
plus que la réserve unique, mais elles se recouvrent largement et
chacune perd les espèces qui ont besoin de plus de \qty{100}{km^2}, donc
moins en pratique. Reliées par des corridors, elles se comportent comme
une seule réserve de \qty{1000}{km^2} pour les espèces qui empruntent
les corridors~: environ un tiers de nouveau, avec l'avantage qu'une
catastrophe locale ne les vide pas toutes.
\textbf{23.} \qty{1000}{km^2}, c'est \qty{31.6}{km} de côté~: intérieur
$(31.6 - 0.6)^2 = \qty{961}{km^2}$, \qty{96}{\%}. \qty{100}{km^2},
c'est \qty{10}{km} de côté~: $(10 - 0.6)^2 = 88$, \qty{88}{\%}.
\textbf{24.} Une lente décroissance du nombre initial vers le nouvel
équilibre $cA^z$, sur des décennies~: les espèces qui s'éteindront sont
d'abord encore présentes, en \omterm{def:b1:populations:population}{populations} trop petites pour persister
--- la dette d'extinction est l'écart entre le compte présent et
l'équilibre, qui reste à payer.
\textbf{25.} $z = 0.25$~; l'île de \qty{100}{km^2} porte \num{38}
espèces à l'équilibre et environ \num{126} une fois reliée par le pont.
\end{solution}
